Seluruh kode program, berkas konfigurasi, model simulasi, serta data pendukung yang digunakan dalam penelitian ini tersedia pada repositori berikut.

\begin{center}
\url{https://github.com/PriagungRA/TA-Autonomous-Robot-Navigation-Fuzzy-Lidar-Thermal}
\end{center}

Repositori tersebut memuat implementasi sistem navigasi robot beroda otonom berbasis fusi sensor LiDAR dan kamera termal, konfigurasi ROS Navigation Stack, implementasi logika fuzzy Mamdani, algoritma \textit{Artificial Potential Field} (APF), berkas simulasi Gazebo, serta skrip pemrosesan sensor yang digunakan selama proses penelitian dan pengujian.

Penyediaan repositori penelitian bertujuan untuk meningkatkan transparansi, reproduktibilitas, dan kemudahan pengembangan lebih lanjut terhadap metode yang diusulkan pada penelitian ini.